\documentclass[10pt,a4paper]{amsart}
\usepackage[margin=19mm]{geometry}
\usepackage{amsmath,amssymb,mathtools}
\usepackage[scr=rsfs,cal=euler]{mathalfa}
\usepackage{xfrac}
\usepackage[unicode,hidelinks,bookmarks=false]{hyperref}
\newcommand{\ie}{i.e.\ }
\newcommand{\cf}{cf.\ }
\newcommand{\resp}{respectively\ }
\newcommand{\Subsection}{\subsection}
\providecommand{\nobreakdash}{\nobreak\hbox{-}}
\setlength{\emergencystretch}{2em}
\multlinegap0pt
\usepackage{mathtools}
\usepackage{amssymb}
\newtheorem{thm}{Theorem}
\newtheorem{lem}{Lemma}
\newcommand{\C}{\mathbb C}
\newcommand{\CA}{\mathcal A}
\newcommand{\R}{\mathbb R}
\title{Superconnections, descent, and monodromy on transversely holomorphic foliations}
\author{Qingyun Zeng}
\date{Source: arXiv:2609.04796v1 (CC BY 4.0)}
\begin{document}
\maketitle

\noindent\emph{Source and licence.} Superconnections, descent, and monodromy on transversely holomorphic foliations. Version arXiv:2609.04796v1 (CC BY 4.0), distributed by arXiv under the Creative Commons Attribution 4.0 International licence, http://creativecommons.org/licenses/by/4.0/, as recorded in the arXiv licence metadata for this exact version and shown on the abstract page https://arxiv.org/abs/2609.04796v1. Full licence notice: \texttt{licenses/arxiv-2609.04796-CC-BY-4.0.txt}.\par\smallskip

\noindent\textit{Editorial scope.} Counted unit: the article's thm thm:suspension-derived-monodromy (source0.tex lines 1357--1370). The article's complete original proof of it is delivered with it (source0.tex lines 1372--1433, one \texttt{proof} environment of 62 lines). No mathematics is written, repaired or re-derived: the statement and the proof are the article's own lines, in the article's own notation, and no retained line is substituted or re-flowed.  Retained uncounted as the source-local context the proof needs: lines 760--786; lines 1268--1279.  Nothing was omitted from any retained span, so the package carries no omission record.  No retained line contains a citation, so the wrapper carries no bibliography and no bibliography entry.  No retained line contains a figure environment, an image-inclusion command or drawing code, so no image is delivered and the package declares no asset dependency.  Editorial formatting only, in addition: an AMS \texttt{amsart} preamble that restates the article's own declarations and macros used by the retained lines, namely the environments \texttt{thm}, \texttt{lem} on separate counters, together with the standard packages those lines need.  The retained environments print in this document as Theorem 1, Lemma 1 in order of appearance, whereas the article prints the same items with the numbering of its own full document.  The complete native source of the article as distributed in the arXiv e-print for this exact version is bundled unchanged as \texttt{sources/209394/source0.tex}.
\begin{thm}[Globally bounded superconnection descent]
\label{thm:eis-cohesive-descent}
Restriction induces a dg quasi-equivalence
\[
\mathcal T_{\mathfrak U}:
\mathcal P_{\CA^\bullet(X)}
\xrightarrow{\ \simeq_{\mathrm{qe}}\ }
\operatorname{Tw}^B_{\mathfrak U}(\mathcal P_{\CA}).
\]
If \(\mathfrak U\) is finite, then
\[
\operatorname{Tw}^B_{\mathfrak U}(\mathcal P_{\CA})
=
\operatorname{Tw}_{\mathfrak U}(\mathcal P_{\CA}),
\]
and hence
\[
\mathcal P_{\CA^\bullet(X)}
\simeq_{\mathrm{qe}}
\operatorname*{holim}_{[p]\in\Delta}
\prod_{(i_0,\ldots,i_p)\in I^{p+1}}
\mathcal P_{\CA^\bullet(U_{i_0\cdots i_p})}.
\]
Empty intersections are omitted.  The same conclusion holds for a locally
finite cover of compact \(X\), since such a cover has only finitely many
nonempty members.
\end{thm}

\begin{lem}[Contractible real factor]
\label{lem:contractible-real-factor}
Let \(I\subset\R\) be an open interval and let
\(p:I\times Z\to Z\) be projection.  Pullback induces a dg
quasi-equivalence
\[
\mathcal P_{\CA^{0,\bullet}(Z)}
\xrightarrow{\ \simeq_{\mathrm{qe}}\ }
\mathcal P_{\Omega^\bullet(I)\widehat\otimes
\CA^{0,\bullet}(Z)}.
\]
\end{lem}

\begin{thm}[Derived monodromy for a holomorphic suspension]
\label{thm:suspension-derived-monodromy}
Restriction to the transversal \(Z\), together with continuation once
around the base circle, induces an equivalence
\[
\boxed{
\mathcal P_{\CA^\bullet(X_\phi)}
\simeq
\mathcal C_Z^{h\phi}
}
\]
in \(\operatorname{Cat}^{\mathrm{perf}}_{\infty,\C}\), equivalently a
Morita equivalence of dg categories.
\end{thm}

\begin{proof}
Cover \(S^1\) by two open arcs \(U_0,U_1\) such that
\[
U_0\cap U_1=J_0\amalg J_1
\]
is the disjoint union of two open intervals.  Choose lifts of the arcs to
\(\R\).  They trivialize the inverse images as \(U_i\times Z\); on one
component of the overlap the transition is the identity of \(Z\), and on
the other it is \(\phi\).

The binary-cover form of
Theorem~\ref{thm:eis-cohesive-descent} gives
\[
\mathcal P_{\CA^\bullet(X_\phi)}
\simeq
\mathcal P_{\CA^\bullet(U_0\times Z)}
\mathop{\times}\limits^h_{
\mathcal P_{\CA^\bullet((J_0\amalg J_1)\times Z)}
}
\mathcal P_{\CA^\bullet(U_1\times Z)}.
\]
All three local categories and the global cohesive category are stable and
idempotent-complete, as noted above.  The displayed dg homotopy pullback
therefore also computes the pullback in
\(\operatorname{Cat}^{\mathrm{perf}}_{\infty,\C}\); this uses closure of
stable idempotent-complete categories under limits, not preservation of
pullbacks by Morita localization.

The category on the disconnected overlap is the product of the categories
on its two components.  Indeed, the central idempotents of the product dga
split every finite projective module and its superconnection; conversely,
two cohesive modules combine over the product dga.

Lemma~\ref{lem:contractible-real-factor} replaces the three vertices of this
cospan by
\[
\mathcal C_Z,\qquad
\mathcal C_Z\times\mathcal C_Z,\qquad
\mathcal C_Z.
\]
Fix one pair of evaluation functors on \(J_0\amalg J_1\).  On the first
leg, evaluation after restriction is smoothly homotopic to evaluation on
\(U_0\) in each component.  On the second leg the same statement holds,
with \(\phi^*\) retained on the component where the trivializations differ.
More explicitly, every evaluation functor is an inverse to the same
constant-extension quasi-equivalence of
Lemma~\ref{lem:contractible-real-factor}.  In an infinity-category the space
of inverses to an equivalence is contractible, so these evaluation functors
are naturally equivalent.  The two resulting squares share the chosen
middle evaluation functor.  Since the indexing cospan has no composable
nonidentity arrows, no further coherence is required.  The resulting
cospan is
\[
\mathcal C_Z
\xrightarrow{(1,1)}
\mathcal C_Z\times\mathcal C_Z
\xleftarrow{(1,\phi^*)}
\mathcal C_Z.
\]
Homotopy limits preserve equivalences of diagrams, and its homotopy
pullback is \(\mathcal C_Z^{h\phi}\) by definition.
\end{proof}

\end{document}
